\documentclass[11pt,letterpaper]{article}

\usepackage[utf8]{inputenc}
\usepackage[T1]{fontenc}
\usepackage[margin=1in]{geometry}
\usepackage{amsmath,amssymb}
\usepackage{booktabs}
\usepackage{enumitem}
\usepackage{xcolor}
\usepackage{titlesec}
\usepackage{abstract}
\usepackage[hidelinks]{hyperref}

\definecolor{njitred}{RGB}{204,0,0}
\definecolor{slate}{RGB}{60,64,72}

\titleformat{\section}{\large\bfseries\color{njitred}}{\thesection}{0.7em}{}
\titleformat{\subsection}{\normalsize\bfseries\color{slate}}{\thesubsection}{0.6em}{}

\newcommand{\code}[1]{\texttt{\small #1}}

\title{\vspace{-2em}\textbf{\color{njitred} GSA Gateway}\\[0.3em]
\large A Local-First, Database-Centric Architecture for Grounded\\
Multi-Platform Conversational Assistance}
\author{Mohammad Dindoost\\\small NJIT Graduate Student Association}
\date{\small Technical Report --- June 2026}

\begin{document}
\maketitle

\begin{abstract}
\noindent
GSA Gateway is a privacy-preserving, locally hosted conversational assistant and
administrative platform serving a university graduate student association across
Discord and Telegram. The system answers free-form questions about events, funding,
policy, and resources using a fully local Retrieval-Augmented Generation (RAG)
pipeline grounded in a curated knowledge base, with no data leaving the host. This
report presents the system's technical architecture and the engineering decisions
behind it. We describe the evolution from a conventional vector-store design to a
unified single-file architecture in which dense vector search (\code{sqlite-vec}),
lexical full-text search (FTS5), relational data, and configuration all live in one
SQLite database; dense and sparse rankings are fused with Reciprocal Rank Fusion. We
present a platform-agnostic \emph{connector} model that fans a single logical message
out to multiple channels with a delivery audit trail, a database-centric data model
(organization hierarchy, versioned knowledge items, posts as a universal output type,
inherited settings), and a methodology for migrating a live production system through
feature-flagged coexistence with instant rollback. We close with a reliability case
study, a qualitative evaluation, and limitations. The recurring theme is a deliberate
bias toward \emph{locality, auditability, and graceful degradation} over scale.
\end{abstract}

\section{Introduction}

Student organizations accumulate institutional knowledge---deadlines, funding rules,
contacts, event logistics---that is chronically hard for members to retrieve. GSA
Gateway addresses this by exposing that knowledge through natural-language chat on the
platforms students already use, while remaining cheap to operate and respectful of
privacy. The design is shaped by four hard requirements:

\begin{enumerate}[leftmargin=1.4em,itemsep=2pt]
  \item \textbf{Grounded answers.} The assistant must answer \emph{only} from a
  curated knowledge base and decline otherwise; fabricated answers are unacceptable in
  an official channel.
  \item \textbf{Privacy and locality.} No student data or query content may leave the
  host. This precludes hosted LLM APIs and external vector databases.
  \item \textbf{Multi-platform reach.} The same assistant and the same broadcasts must
  serve Discord and Telegram, with room for future channels.
  \item \textbf{Low operational burden.} The system runs on a single always-on machine
  and is administered by a non-specialist officer.
\end{enumerate}

\noindent
These constraints drive most of what follows. The technical contributions of the
system are: (i) a fully local RAG stack (embeddings and generation via Ollama) with
principled grounding and degradation; (ii) a \emph{single-database} hybrid retrieval
engine that replaces a separate vector store by combining \code{sqlite-vec} KNN and
FTS5 BM25 under Reciprocal Rank Fusion; (iii) a connector abstraction decoupling
message content from platform-specific delivery, with a persistent delivery log;
(iv) a database-centric, multi-tenant data model with versioned knowledge and
inherited settings; and (v) a coexistence strategy that let a second-generation
architecture be introduced into a running system behind feature flags, with drop-in
interface adapters and instant rollback.

\section{System Architecture}

GSA Gateway runs as a single bot process (with a companion Telegram poller) on one
host. Architecturally it is layered:

\begin{itemize}[leftmargin=1.4em,itemsep=2pt]
  \item \textbf{Interaction layer.} Slash commands and an \code{on\_message} handler
  for free-form chat; an intent detector routes greetings, thanks, and trivial
  queries away from the expensive RAG path.
  \item \textbf{Retrieval layer.} Converts a question into grounding context (the
  focus of \S\ref{sec:rag} and \S\ref{sec:hybrid}).
  \item \textbf{Generation layer.} A local LLM produces an answer constrained to the
  retrieved context, or the system degrades to lexical search.
  \item \textbf{Data layer.} A SQLite database holds operational records and, in the
  second-generation design, the entire knowledge and configuration substrate.
  \item \textbf{Delivery layer.} A connector registry transmits outbound messages to
  every targeted platform (\S\ref{sec:connectors}).
\end{itemize}

A defining characteristic is that the system exists in two coexisting generations. The
first (\code{v1}) is the deployed application built around a dedicated vector database
(ChromaDB). The second (\code{v2}) is a database-first redesign that consolidates
storage and adds hybrid retrieval; it is wired \emph{into} the running v1 process
behind feature flags rather than deployed as a separate service. \S\ref{sec:method}
explains why this coexistence was the safest migration path.

\section{Grounded Retrieval-Augmented Generation}
\label{sec:rag}

\subsection{Offline indexing}
Knowledge-base documents are split into chunks of at most $\approx\!350$ tokens at
semantic boundaries, balancing two competing pressures: chunks must be small enough
that a retrieved unit is densely relevant, yet large enough to carry self-contained
meaning. Each chunk is embedded with \code{nomic-embed-text} (768 dimensions) served
locally by Ollama. The model is \emph{asymmetric}: documents are embedded with a
\code{search\_document:} prefix and queries with \code{search\_query:}. Using the
wrong prefix silently degrades recall, so the distinction is enforced at the API
boundary. Embeddings are $L_2$-normalized at write time so that the default Euclidean
index ranks identically to cosine similarity:
\[
\lVert a - b\rVert_2^2 = 2\,(1 - \cos(a,b)) \quad\text{for } \lVert a\rVert=\lVert b\rVert=1 .
\]

\subsection{Online answering}
A query is embedded, matched against the index by approximate nearest neighbor,
reranked, and the top chunks become grounding context for the LLM
(\code{llama3.1:8b}). The system prompt instructs the model to answer strictly from
the provided context and to refuse when the answer is absent---the mechanism that
keeps an official channel from emitting confident falsehoods. Per-user conversation
memory (a short, time-bounded window) supports follow-up questions without unbounded
context growth.

\subsection{Graceful degradation}
If the LLM is unavailable, the assistant does not fail: it falls back to fuzzy lexical
matching (\code{rapidfuzz}) over the knowledge base, returning the best matching entry
below a confidence threshold rather than a generated answer. This makes the heavy
generation component optional---valuable on a resource-constrained host, where it can
be disabled deliberately to free memory.

\section{Hybrid Retrieval in a Single Database}
\label{sec:hybrid}

\subsection{Motivation: collapsing the storage tier}
The first-generation design used three stores: ChromaDB for vectors, SQLite for
records, and in-memory fuzzy search. This works but multiplies operational surface:
separate persistence, separate backup semantics, and no transactional consistency
between an item and its embedding. The second generation collapses all of it into one
SQLite file using two extensions: \code{sqlite-vec} provides a \code{vec0} virtual
table for native KNN over \code{FLOAT[768]} vectors, and FTS5 provides BM25-ranked
full-text search. Triggers keep the FTS index synchronized with the canonical
\code{knowledge\_items} table. The payoff is a single portable artifact with atomic
writes, one backup, and no cross-store drift.

\subsection{Fusing dense and sparse signals}
Dense (semantic) retrieval captures paraphrase and conceptual similarity but can miss
exact tokens; sparse (lexical) retrieval excels at names, acronyms, and rare terms but
misses synonymy. The system runs both and fuses their \emph{rankings} with Reciprocal
Rank Fusion (RRF):
\[
\mathrm{score}(d) \;=\; \sum_{i \in \{\text{dense},\,\text{sparse}\}}
\frac{1}{k + \mathrm{rank}_i(d)}, \qquad k = 60 .
\]
RRF combines result lists by rank rather than by raw score, which sidesteps the
incompatibility between cosine similarities and BM25 magnitudes and is robust without
per-query tuning. A small type-dependent boost is then applied (e.g.\ favoring
\code{contact} and \code{event} items) to reflect that certain question classes have
canonical answer types.

\subsection{A production lesson: decoupling the candidate pool}
An early failure illustrates a subtle retrieval pitfall. The candidate pool fed to
fusion was initially sized as a multiple of the requested result count
($\text{pool}=\max(c\cdot\text{limit},\,15)$). For a query whose correct answer ranked
highly in the \emph{sparse} list but only moderately in the \emph{dense} list, a small
pool truncated the answer out of the dense candidates before fusion could rescue it,
and a less relevant item was returned. The fix was to decouple the pool from the
result count---using a fixed, generous candidate pool (e.g.\ 40) independent of how
many results the caller ultimately wants. The general principle: \emph{fusion quality
is bounded by candidate recall}, so the pool must be sized for recall, not for the
final cut.

\section{Data Model}
\label{sec:data}

The second-generation schema is deliberately database-centric: structure, content, and
configuration are all rows, not code or files.

\begin{itemize}[leftmargin=1.4em,itemsep=3pt]
  \item \textbf{Organization hierarchy.} An \code{organizations} table forms a tree via
  \code{parent\_id} (university $\to$ college $\to$ department $\to$ club, etc.),
  traversed with recursive common table expressions. This makes the platform
  multi-tenant: a single deployment can host many sub-organizations.
  \item \textbf{Versioned knowledge.} \code{knowledge\_items} are typed and versioned:
  a \code{root\_id} groups all versions of an item, \code{parent\_id} links to the
  prior version, and \code{is\_active} marks the current one. Edits append a new
  version rather than mutating in place, preserving history and enabling restore.
  \item \textbf{Posts as universal output.} Every outbound message---one-off, recurring
  template, or event announcement---is a row in a single \code{posts} model, decoupling
  \emph{what} is said from \emph{where} it is delivered.
  \item \textbf{Inherited settings.} A \code{settings} table stores configuration per
  organization; lookups walk the ancestor chain (self $\to$ parent $\to$ root), so a
  value set at the root is inherited by all descendants unless overridden. The same
  mechanism stores controlled \emph{vocabularies} (valid organization types, knowledge
  types, contact roles), letting each deployment define its taxonomy as data.
\end{itemize}

\noindent
Tables use SQLite's \code{STRICT} mode for type enforcement, and all timestamps are
stored as canonical UTC and rendered in the organization's local timezone at the edge,
avoiding a class of off-by-hours scheduling bugs.

\section{Multi-Platform Delivery: the Connector Pattern}
\label{sec:connectors}

Outbound messaging is built around a single abstraction. A \code{Post} describes
\emph{what} to send (content, optional signature, target platforms and channels). A
\code{BaseConnector} subclass knows how to send for \emph{one} platform
(\code{send\_text}, \code{send\_media}, \code{send\_interactive}). A
\code{ConnectorRegistry} fans a post out to all targeted connectors in parallel and---
critically---\emph{never raises}: each delivery is wrapped, failures become recorded
failure results, and every attempt is written to a \code{post\_deliveries} audit table.

This yields three properties. First, \textbf{extensibility}: adding a platform is one
new subclass and a registration, with no change to callers. Second, \textbf{isolation}:
a misbehaving platform cannot sink a multi-channel broadcast. Third,
\textbf{auditability}: \code{post\_deliveries} is a durable record of what was sent
where, the substrate for a ``single source of truth'' delivery policy. A scheduler
materializes due recurring templates and event reminders into concrete posts and hands
them to the registry, all on UTC time.

\section{Case Study: Fault-Tolerant Live Event Tracking}

A live sports tracker exercises the system differently from Q\&A: it polls an external
API on a fixed cadence, diffs successive snapshots of match state to emit
\emph{events} (kickoff, goal, half-time, full-time), deduplicates via per-match state,
and publishes each event through the connector registry. Two reliability lessons
emerged.

First, a long-lived HTTP client session that was reused across polls suffered a
$100\%$ failure rate in production while a one-shot request to the same endpoint
succeeded: the pooled TCP connection went stale between polls and every reuse was reset.
The remedy was to use a fresh client session per poll with bounded retries. Second, the
upstream API itself proved intermittently unavailable (roughly one connection in six
was dropped even for fresh sessions). Because the tracker polls on a short cadence and
each event need only be reported once within a small window, retries combined with the
poll interval yield effectively complete coverage despite a flaky dependency. The
takeaway generalizes: for periodic integrations, prefer \emph{stateless} requests and
design for an unreliable upstream rather than assuming a healthy long-lived connection.

\section{Engineering Methodology}
\label{sec:method}

\subsection{Feature-flagged coexistence and instant rollback}
The most consequential decision was \emph{not} to deploy the second generation as a
replacement. Instead, v2 components are imported into the running v1 process and gated
by environment flags. A drop-in adapter (a ``shim'') presents the exact interface the
v1 code expects, so the new hybrid retriever can substitute for the old vector-store
retriever transparently. Any flag can be flipped off and the process restarted to
return to v1 behavior in seconds. This converts a high-risk rewrite into a sequence of
independently reversible steps---essential when the system is a live service with a
fixed external deadline.

\subsection{Safe data migration}
Schema and data migrations are idempotent and guarded: they run first against a copy of
the live database (a dry run), and any live run is preceded by an automatic, timestamped
backup. Administrative writes from the dashboard are likewise idempotent
(\code{INSERT OR REPLACE}), so re-applying a change is harmless.

\subsection{A dual-mode administrative interface}
Administration is intentionally dependency-free. The dashboard is static
HTML/CSS/JavaScript that reads a SQLite database in the browser via WebAssembly
(\code{sql.js}). It operates in two modes. In \emph{file mode} it reads a database copy
and emits changes as SQL patch scripts to be applied out of band---never writing the
live file directly. In \emph{server mode} a tiny standard-library HTTP server on the
host (bound to localhost only, reached over an SSH tunnel) serves the dashboard and
mediates reads and writes against the live database, returning WAL-consistent
snapshots. The same UI thus supports both a zero-trust ``inspect and export'' workflow
and a fully interactive one, without shipping a heavyweight web framework.

\subsection{Single-source-of-truth governance}
Because any process holding platform credentials can post directly, ``only allow posts
that originate in the database'' cannot be enforced at the platform boundary. The
practical model is constructive: route all sends through the connector registry (which
logs them), keep credentials in a single process, retire ad-hoc senders, and reconcile
channel activity against \code{post\_deliveries} to detect anything out of band.

\section{Evaluation}

Evaluation to date is qualitative and operational rather than a formal benchmark, which
we state plainly as a limitation. Retrieval behavior was validated against a suite of
representative queries spanning funding, contacts, and cross-domain questions; the
hybrid-with-fixed-pool design corrected concrete failures of the earlier vector-only
retriever (notably surfacing the correct office for finance questions). Correctness of
the non-retrieval machinery---schema, hybrid ranking, connectors, scheduling, the
administrative server, and event detection/deduplication---is covered by an automated
unit-test suite, alongside the inherited first-generation tests. Operationally, the
system has run as an always-on service, and the connector audit trail provides
end-to-end confirmation of multi-platform delivery.

\section{Limitations and Future Work}

The deployment is single-host and vertically scaled; horizontal scale is out of scope
and arguably unnecessary at this size. Answer quality and latency are bounded by a
local 8-billion-parameter model. Retrieval is validated qualitatively; a labeled
relevance benchmark would enable quantitative tuning of the fusion constant and type
boosts. The taxonomy vocabularies guide but do not constrain values at the database
level. Finally, the two generations still coexist: completing the migration entails
re-implementing the remaining interaction surface (the full command set and ancillary
schedulers) on the v2 substrate so the first generation can be retired, and unifying
rich, platform-specific message formatting within the connector model.

\section{Conclusion}

GSA Gateway shows that a small, local-first system can deliver grounded,
multi-platform conversational assistance without external dependencies. Its more
transferable ideas are architectural: consolidating dense vectors, lexical search, and
relational data into a single SQLite file with rank-fusion retrieval; separating
message content from platform delivery behind an auditable connector registry; and
migrating a live system safely through feature-flagged coexistence with drop-in
adapters and instant rollback. Throughout, the design trades scale for locality,
auditability, and graceful degradation---an appropriate and, we argue, principled
choice for community-scale software.

\end{document}
